% =====================================================================
\chapter{El producto de matrices}
\label{tema:4}
% =====================================================================

El producto de matrices se define de modo que represente la composición de aplicaciones lineales. De esa definición se siguen sus propiedades (asociativa, distributivas, comportamiento con la traspuesta) y su principal diferencia con el producto de números: no es conmutativo. El tema sigue la sección~2.4 de Strang~\cite{strang} y la sesión «Multiplication and Inverse Matrices» de 18.06SC~\cite{mit1806sc}; la composición y el coste computacional se tratan como en~\cite[§10.1 y §10.2]{vmls}.

\section{Definición y composición}

\begin{definicion}[Producto de matrices]\label{def:producto}
Sean $A\in\R^{m\times n}$ y $B\in\R^{n\times p}$. El producto $AB\in\R^{m\times p}$ es la matriz de entradas
\[
(AB)_{ij}=\sum_{k=1}^{n}A_{ik}B_{kj},\qquad 1\le i\le m,\ 1\le j\le p.
\]
La entrada $(i,j)$ es el producto escalar de la fila $i$ de $A$ por la columna $j$ de $B$. El producto solo está definido si el número de columnas de $A$ coincide con el número de filas de $B$: $(m\times n)(n\times p)=m\times p$.
\end{definicion}

\begin{teorema}[El producto es la composición]\label{teo:composicion}
Sean $A\in\R^{m\times n}$ y $B\in\R^{n\times p}$. Para todo $x\in\R^{p}$, $(AB)x=A(Bx)$. Es decir, la matriz de la aplicación «aplicar primero $B$ y después $A$» es $AB$.
\end{teorema}

\begin{proof}
La componente $i$ de $A(Bx)$ es
\[
\sum_{k=1}^{n}A_{ik}(Bx)_k=\sum_{k=1}^{n}A_{ik}\sum_{j=1}^{p}B_{kj}x_j=\sum_{j=1}^{p}\Bigl(\sum_{k=1}^{n}A_{ik}B_{kj}\Bigr)x_j=\sum_{j=1}^{p}(AB)_{ij}x_j,
\]
que es la componente $i$ de $(AB)x$. Se ha usado solo que las sumas finitas pueden reordenarse.
\end{proof}

El orden de escritura es el de la composición de funciones, $(A\circ B)(x)=A(B(x))$: en $AB$ actúa primero la matriz de la derecha.

\begin{ejemplo}\label{ej:intercambio}
Sean $A=\begin{pmatrix}1&2\\3&4\end{pmatrix}$ y $B=\begin{pmatrix}0&1\\1&0\end{pmatrix}$. Entonces
\[
AB=\begin{pmatrix}2&1\\4&3\end{pmatrix},\qquad BA=\begin{pmatrix}3&4\\1&2\end{pmatrix}.
\]
Multiplicar por $B$ a la derecha intercambia las columnas de $A$; multiplicar por $B$ a la izquierda intercambia sus filas.
\end{ejemplo}

\section{Cuatro formas de calcular un producto}

La \cref{def:producto} calcula el producto entrada a entrada. Strang subraya otras tres formas equivalentes~\cite[§2.4]{strang}, cada una útil en contextos distintos.

\begin{proposicion}\label{prop:cuatro}
Sean $A\in\R^{m\times n}$ con columnas $a_1,\dots,a_n$ y filas $r_1^{\top},\dots,r_m^{\top}$, y $B\in\R^{n\times p}$ con columnas $b_1,\dots,b_p$ y filas $s_1^{\top},\dots,s_n^{\top}$. Entonces:
\begin{enumerate}
  \item (por columnas) la columna $j$ de $AB$ es $Ab_j$;
  \item (por filas) la fila $i$ de $AB$ es $r_i^{\top}B$;
  \item (columnas por filas) $AB=\sum_{k=1}^{n}a_k\,s_k^{\top}$, suma de $n$ matrices $m\times p$.
\end{enumerate}
\end{proposicion}

\begin{proof}
(i) La componente $i$ de $Ab_j$ es $\sum_kA_{ik}B_{kj}=(AB)_{ij}$. (ii) La componente $j$ de $r_i^{\top}B$ es $\sum_kA_{ik}B_{kj}$. (iii) La entrada $(i,j)$ de la matriz $a_ks_k^{\top}$ es $A_{ik}B_{kj}$; sumando en $k$ se obtiene $(AB)_{ij}$.
\end{proof}

Una matriz de la forma $a\,s^{\top}$, con $a\in\R^{m}$ y $s\in\R^{p}$, se llama \emph{producto exterior}; todas sus filas son múltiplos de $s^{\top}$. La forma (iii) descompone un producto de matrices en suma de productos exteriores y es la base de la descomposición en valores singulares y de las aproximaciones de bajo rango del módulo~3. En el \cref{ej:intercambio}, $a_1s_1^{\top}+a_2s_2^{\top}=\begin{pmatrix}0&1\\0&3\end{pmatrix}+\begin{pmatrix}2&0\\4&0\end{pmatrix}=AB$.

\section{Propiedades algebraicas}

\begin{proposicion}\label{prop:algebra}
Siempre que los tamaños permitan formar los productos:
\begin{enumerate}
  \item (asociativa) $(AB)C=A(BC)$;
  \item (distributivas) $A(B+C)=AB+AC$ y $(B+C)A=BA+CA$;
  \item (identidad) $I_mA=A=AI_n$ si $A\in\R^{m\times n}$;
  \item (traspuesta) $(AB)^{\top}=B^{\top}A^{\top}$.
\end{enumerate}
\end{proposicion}

\begin{proof}
(i) Por el \cref{teo:composicion}, para todo $x$, $((AB)C)x=(AB)(Cx)=A(B(Cx))$ y $(A(BC))x=A((BC)x)=A(B(Cx))$. Dos matrices que actúan igual sobre todo vector son iguales (basta tomar $x=e_j$ para ver que tienen las mismas columnas). (ii) $\sum_kA_{ik}(B_{kj}+C_{kj})=\sum_kA_{ik}B_{kj}+\sum_kA_{ik}C_{kj}$, y análogamente por la derecha. (iii) Solo $k=i$ contribuye en $\sum_k(I_m)_{ik}A_{kj}$. (iv) $((AB)^{\top})_{ij}=(AB)_{ji}=\sum_kA_{jk}B_{ki}=\sum_k(B^{\top})_{ik}(A^{\top})_{kj}=(B^{\top}A^{\top})_{ij}$.
\end{proof}

La distributiva por la izquierda y la distributiva por la derecha son propiedades distintas: no puede pasarse de una a otra cambiando el lado (ejercicio~\ref{eje:4-distributiva}), porque el producto no es conmutativo.

\begin{observacion}[No conmutatividad]\label{obs:no-conmuta}
En general $AB\neq BA$, incluso si ambos productos están definidos y son del mismo tamaño. Sean $R=\begin{pmatrix}0&-1\\1&0\end{pmatrix}$ (rotación de $90^\circ$), $S=\begin{pmatrix}2&0\\0&0{,}5\end{pmatrix}$ (escalado) y $x=(3,1)$. Escalar y después rotar da $R(Sx)=R(6;\,0{,}5)=(-0{,}5;\,6)$; rotar y después escalar da $S(Rx)=S(-1,3)=(-2;\,1{,}5)$. El orden de las transformaciones importa.
\end{observacion}

\begin{definicion}[Matriz ortogonal]\label{def:ortogonal}
Una matriz $Q\in\R^{n\times n}$ es \emph{ortogonal} si $Q^{\top}Q=I_n$, es decir, si sus columnas son unitarias y ortogonales dos a dos.
\end{definicion}

\begin{proposicion}\label{prop:ortogonal}
Si $Q$ es ortogonal, entonces $(Qu)\cdot(Qv)=u\cdot v$ para todos $u,v\in\R^{n}$. En consecuencia, $Q$ conserva normas, distancias y ángulos.
\end{proposicion}

\begin{proof}
Por la \cref{prop:algebra}(iv), $(Qu)^{\top}(Qv)=u^{\top}Q^{\top}Qv=u^{\top}I_nv=u^{\top}v$. Las normas, las distancias y los ángulos se expresan mediante productos escalares.
\end{proof}

Las rotaciones $R_\theta$ son ortogonales (ejercicio~\ref{eje:4-ortogonal}), lo que generaliza la \cref{prop:rotacion}. Las codificaciones posicionales rotatorias (RoPE) del módulo~14 aplican precisamente rotaciones de este tipo a las consultas y las claves.

\section{Coste y matriz de Gram}

\begin{proposicion}[Coste del producto]\label{prop:coste}
Calcular $AB$ con $A\in\R^{m\times n}$ y $B\in\R^{n\times p}$ mediante la \cref{def:producto} requiere $mnp$ multiplicaciones y $mp(n-1)$ sumas.
\end{proposicion}

\begin{proof}
Cada una de las $mp$ entradas es un producto escalar en $\R^{n}$, que requiere $n$ multiplicaciones y $n-1$ sumas.
\end{proof}

Aunque $(AB)C=A(BC)$, el coste de los dos órdenes puede ser muy distinto (ejercicio~\ref{eje:4-coste}). Elegir bien el orden de los paréntesis es una optimización habitual en el cálculo con tensores.

\begin{definicion}[Matriz de Gram]\label{def:gram}
Sea $X\in\R^{n\times d}$ la matriz cuyas filas son los vectores $x_1^{\top},\dots,x_n^{\top}$. Su \emph{matriz de Gram} es $G=XX^{\top}\in\R^{n\times n}$, con $G_{ij}=x_i\cdot x_j$.
\end{definicion}

La matriz de Gram es simétrica, porque $G^{\top}=(XX^{\top})^{\top}=XX^{\top}$, y su diagonal contiene las normas al cuadrado, $G_{ii}=\norm{x_i}^{2}$. Un único producto de matrices calcula los $n^{2}$ productos escalares entre $n$ vectores.

\begin{ejemplo}\label{ej:gram}
Con los embeddings del \cref{ej:juguete}, sea $X$ la matriz $3\times 4$ de filas «rey», «reina» y «manzana». Entonces
\[
XX^{\top}=\begin{pmatrix}9&7&0\\7&9&0\\0&0&4\end{pmatrix}.
\]
La entrada $(1,2)$ es $\text{rey}\cdot\text{reina}=7$, y dividiendo por $\sqrt{G_{11}G_{22}}=9$ se recupera la similitud coseno $7/9$.
\end{ejemplo}

\begin{observacion}
En la atención, las consultas y las claves de los $n$ tokens de una secuencia se apilan como filas de matrices $Q,K\in\R^{n\times d_k}$, y todas las puntuaciones se calculan con un único producto $QK^{\top}\in\R^{n\times n}$, que generaliza la matriz de Gram. Las GPU están diseñadas para ejecutar este tipo de productos con gran rapidez (módulo~17).
\end{observacion}

\begin{lecturas}
\item Strang~\cite{strang}, §2.4, reglas de las operaciones con matrices, y §2.7, traspuestas.
\item MIT 18.06SC~\cite{mit1806sc}, unidad~I, sesión «Multiplication and Inverse Matrices».
\item Deisenroth, Faisal y Ong~\cite{mml}, §2.2 «Matrices» y §3.9 «Rotations».
\item Boyd y Vandenberghe~\cite{vmls}, §10.1 «Matrix-matrix multiplication» y §10.2 «Composition of linear functions».
\end{lecturas}

\section*{Ejercicios del tema 4}
\addcontentsline{toc}{section}{Ejercicios del tema 4}

\begin{ejercicio}{1}{eje:4-formas}
(a) ¿Qué tamaño tiene el producto de una matriz $3\times 4$ por otra $4\times 2$? ¿Puede multiplicarse una matriz $2\times 3$ por otra $2\times 3$? (b) Para $u,v\in\R^{3}$, ¿qué tamaños tienen $u^{\top}v$ y $uv^{\top}$? (c) Justifique si es verdadero o falso que $AB=BA$ para todo par de matrices cuadradas del mismo orden.
\end{ejercicio}

\begin{ejercicio}{2}{eje:4-rs}
Con $R$ y $S$ de la \cref{obs:no-conmuta}, calcule $RS$ y $SR$, aplíquelas a $x=(3,1)$ y compruebe que se obtiene $R(Sx)$ y $S(Rx)$, respectivamente.
\end{ejercicio}

\begin{ejercicio}{2}{eje:4-ab-ba}
Sean $A=\begin{pmatrix}1&0&2\\0&1&-1\end{pmatrix}$ y $B=\begin{pmatrix}1&2\\0&1\\3&0\end{pmatrix}$. Calcule $AB$ y $BA$ e indique sus tamaños. ¿Puede calcularse $AA$?
\end{ejercicio}

\begin{ejercicio}{2}{eje:4-distributiva}
Sean $A=\begin{pmatrix}2&1\\1&3\end{pmatrix}$, $B=\begin{pmatrix}1&0\\0&0\end{pmatrix}$ y $C=\begin{pmatrix}0&0\\2&1\end{pmatrix}$. Calcule $AB+AC$ y $A(B+C)$ y compárelas. Calcule también $BA+CA$. ¿Coincide con $A(B+C)$?
\fuente{adaptado de MIT 18.06SC~\cite{mit1806sc}, unidad~I, sesión «Multiplication and Inverse Matrices», problema~3.1 (matrices cambiadas).}
\end{ejercicio}

\begin{ejercicio}{2}{eje:4-coste}
Sean $A\in\R^{3\times 5}$, $B\in\R^{5\times 2}$ y $C\in\R^{2\times 4}$. (a) Con la \cref{prop:coste}, calcule el número de multiplicaciones de $(AB)C$ y de $A(BC)$. (b) Para $u,v,w\in\R^{N}$, compare el coste de $(u^{\top}v)\,w^{\top}$ y de $u^{\top}(vw^{\top})$, y evalúelo para $N=4096$.
\fuente{adaptado de Strang~\cite{strang}, §2.4, problema~36, propuesto en el curso 18.06 del MIT~\cite{mit1806}, primavera de 2010, hoja de problemas~1 (dimensiones cambiadas).}
\end{ejercicio}

\begin{ejercicio}{2}{eje:4-proyeccion}
Sea $A\in\R^{3\times 2}$ la matriz formada por las dos primeras columnas de $I_3$, y sea $b=(4,-1,2)$. (a) Como las columnas de $A$ son unitarias y ortogonales, la matriz de proyección sobre el plano que generan es $P=AA^{\top}$. Calcule $P$ y $Pb$. (b) Compruebe que $P^{2}=P$ y que $(I-P)^{2}=I-P$. ¿Sobre qué proyecta $I-P$? (c) Escriba la proyección sobre la recta de $a=(1,1,1)$ de la \cref{prop:proyeccion} como $P_a\,b$ con $P_a=aa^{\top}/(a^{\top}a)$ y compruebe el resultado del \cref{ej:proyeccion}.
\fuente{adaptado de MIT 18.06SC~\cite{mit1806sc}, unidad~II, sesión «Projections onto Subspaces», problemas~15.1 y~15.2 (dimensión y vector cambiados).}
\end{ejercicio}

\begin{ejercicio}{3}{eje:4-ortogonal}
(a) Demuestre que $R_\theta^{\top}R_\theta=I_2$ para todo $\theta$. (b) Demuestre que $R_\theta R_\varphi=R_{\theta+\varphi}$ y deduzca que dos rotaciones del plano siempre conmutan. (c) Demuestre que el producto de dos matrices ortogonales del mismo orden es ortogonal. (d) Contraste (b) con la \cref{obs:no-conmuta}: ¿qué propiedad de $S$ impide que conmute con $R$?
\end{ejercicio}
